\chapter{L'apparato digerente: stomaco, fegato, pancreas e intestino}

% ---------------------------------------------------------------------------
\section{Lo stomaco}

Lo stomaco è un organo cavo, intraperitoneale, situato nello spazio
sopra-mesocolico della cavità addominale, dove occupa l'ipocondrio
sinistro e l'epigastrio. Ha la forma di un sacco, più ampio nella parte
superiore e schiacciato in senso anteroposteriore. Anteriormente è in
rapporto con il diaframma, il fegato e la parete addominale;
posteriormente con il pancreas e la milza. Tra la parte posteriore dello
stomaco e il pancreas si trova uno spazio rivestito da peritoneo, la
\textbf{borsa omentale}. I margini destro e sinistro dello stomaco sono
curvi e formano rispettivamente la piccola e la grande curvatura.
L'irrorazione arteriosa dello stomaco deriva dai rami del tripode celiaco,
in particolare dalle arterie gastriche e gastroepiploiche, che decorrono
lungo la piccola e la grande curvatura, e da rami dell'arteria splenica.

% ---------------------------------------------------------------------------
\subsection{Le porzioni dello stomaco}

Nello stomaco si distinguono quattro porzioni: il \textbf{cardias}, la
regione circostante lo sbocco dell'esofago; il \textbf{fondo}, a forma di
cupola, la porzione più alta, in contatto con la concavità
diaframmatica; il \textbf{corpo}, la porzione più ampia, di forma
vagamente cilindrica orientata in senso verticale; la \textbf{parte
pilorica} (comprendente antro e canale pilorico), di forma conica e
direzione obliqua verso destra. Le funzioni dello stomaco sono
l'accumulo del cibo ingerito, la sua trasformazione meccanica e la
digestione chimica per azione di acidi ed enzimi.

% ---------------------------------------------------------------------------
\subsection{Parete gastrica e ghiandole gastriche}

La parete dello stomaco presenta una tonaca mucosa il cui epitelio è
formato da abbondanti cellule colonnari a secrezione mucosa, disposte su
un unico strato, e da cellule provviste di microvilli (\textit{brush
cell}), capaci di analizzare il contenuto dello stomaco. Nella lamina
propria sono ospitate le ghiandole gastriche, la cui parete è formata da
diversi elementi cellulari che producono il succo gastrico, alcuni
ormoni e muco fortemente alcalino; le ghiandole gastriche propriamente
dette si trovano nel corpo e nel fondo e sono formate da cellule a
secrezione mucosa (cellule del colletto), cellule principali, cellule
parietali, cellule staminali e cellule neuroendocrine. Le cellule
colonnari producono un secreto mucoso alcalino che difende lo stomaco
dall'attacco dei suoi stessi succhi acidi.

La tonaca mucosa si solleva in pliche e contiene il plesso sottomucoso;
le pliche sono più lisce in corrispondenza del fondo e del cardias,
mentre risultano più rilevate e regolari in prossimità dell'antro e del
canale pilorico. La tonaca muscolare è organizzata in tre strati (uno
profondo a fibre oblique, uno intermedio a fibre circolari, uno
superficiale a fibre longitudinali) e contiene il plesso mienterico.
La giunzione gastroesofagea è costituita dalle formazioni anatomiche che
aiutano a mantenere una porzione dell'esofago e lo stomaco nella cavità
addominale.

% ---------------------------------------------------------------------------
\subsection{Le cellule delle ghiandole gastriche}

Le \textbf{cellule mucose} secernono muco. Le \textbf{cellule
parietali} secernono acido cloridrico, che abbassa il pH del succo
gastrico e uccide i microrganismi, degradando le pareti cellulari e i
tessuti connettivi presenti nel cibo; secernono inoltre il fattore
intrinseco, necessario per l'assorbimento della vitamina B12, e la
grelina, che regola l'appetito. Le \textbf{cellule principali} secernono
il pepsinogeno, precursore della pepsina, enzima che digerisce le
proteine, e la lipasi gastrica, che digerisce i grassi del latte. Le
\textbf{cellule enterocromaffini} secernono ormoni e sostanze che
regolano la digestione; tra queste le cellule G secernono gastrina, che
stimola le ghiandole gastriche a secernere acidi ed enzimi e intensifica
il rimescolamento promuovendo l'attività della muscolatura liscia della
parete gastrica.

% ---------------------------------------------------------------------------
\section{Il pancreas}

Il pancreas è una ghiandola tubulo-acinosa composta, extramurale,
associata al canale alimentare, situata nella cavità addominale
all'altezza delle prime due vertebre lombari. Ha una forma a martello,
con la testa in rapporto stretto con la "C" duodenale, un collo o istmo,
un corpo e una coda che si portano in alto e a sinistra. Il parenchima
pancreatico comprende una componente ghiandolare esocrina, organizzata
come ghiandola tubuloacinosa composta a secrezione sierosa simile a
quella della ghiandola parotide, e una componente endocrina, formata da
piccoli agglomerati di tessuto endocrino detti isole pancreatiche, che
producono ormoni (tra cui insulina e glucagone, immessi direttamente nel
circolo sanguigno) responsabili della regolazione del metabolismo dei
carboidrati. La componente esocrina è fondamentale nella digestione,
poiché è responsabile della produzione del succo pancreatico.

% ---------------------------------------------------------------------------
\subsection{Istologia e secrezione pancreatica}

L'unità funzionale della componente esocrina è l'\textbf{acino
pancreatico}, formato da cellule acinose pancreatiche, che secernono gli
enzimi digestivi e presentano un ricco reticolo endoplasmatico rugoso, un
sviluppato complesso di Golgi e numerosi granuli di zimogeno, e da
cellule centroacinose. La secrezione, raccolta da canalicoli
intracellulari, defluisce nei dotti intercalari e quindi nei dotti
intralobulari, fino a confluire nel dotto pancreatico principale, che si
unisce al dotto coledoco.

Gli acini secernono il succo pancreatico nel duodeno. Questo succo è
caratterizzato dalla presenza di acqua, ioni ed enzimi digestivi, tra i
quali: le lipasi, che agiscono sui lipidi; le carboidrasi, che agiscono
su zuccheri e amidi; le nucleasi, che agiscono sugli acidi nucleici; gli
enzimi proteolitici, tra cui le proteasi, che scindono le grosse
molecole proteiche, e le peptidasi, che trasformano le piccole catene
peptidiche in aminoacidi semplici. Il succo pancreatico viene prodotto
durante la fase cefalica della digestione, in ragione di circa 1 litro al
giorno, e ha un pH compreso tra 7,5 e 8,8, essendo ricco di ioni
bicarbonato secreti dalle cellule dei dotti escretori degli acini
pancreatici.

% ---------------------------------------------------------------------------
\section{Il fegato}

Il fegato è un organo parenchimatoso disposto nella parte superiore
della cavità addominale, nello spazio sopramesocolico; occupa
l'ipocondrio destro, l'epigastrio e una parte dell'ipocondrio sinistro.
È un organo intraperitoneale, in quanto completamente rivestito dal
peritoneo tranne che per una piccola porzione. Ha forma ovoide e riceve
una vascolarizzazione in entrata sia arteriosa sia venosa, attraverso
l'ilo del fegato, disposto sulla sua faccia inferiore: il sangue
arterioso giunge tramite l'arteria epatica propria, quello venoso tramite
la vena porta. Il parenchima epatico è organizzato in lobuli epatici,
che in sezione trasversale appaiono come aree poligonali al cui interno
gli epatociti sono organizzati in cordoni. Il fegato è mantenuto in sede
da diversi legamenti: il legamento falciforme e il legamento rotondo del
fegato, il piccolo omento (comprendente il legamento epatoduodenale e il
legamento epatogastrico) e il legamento gastrocolico, che lo pongono in
rapporto con lo stomaco, la milza e il diaframma.

Il fegato è la voluminosa ghiandola extramurale annessa al sistema
digerente e rappresenta l'organo più voluminoso del corpo umano: pesa
circa 1,5 kg, cui si aggiungono circa 500 g di sangue che vi circolano,
rappresentando complessivamente circa il 2\% del peso corporeo.

% ---------------------------------------------------------------------------
\subsection{Funzioni del fegato}

Tra le funzioni digestive e metaboliche del fegato si annoverano: la
sintesi di somatomedine; la sintesi e secrezione di bile; l'accumulo di
glicogeno e di lipidi di riserva; il controllo delle concentrazioni
ematiche di glucosio, aminoacidi e acidi grassi; la sintesi e
l'interconversione di sostanze nutritizie, ad esempio la transaminazione
degli aminoacidi o la conversione di carboidrati in lipidi; la sintesi e
il rilascio di colesterolo legato a proteine vettrici; l'inattivazione di
tossine; l'accumulo di riserve di ferro; l'accumulo di vitamine
liposolubili.

Tra le altre funzioni fondamentali del fegato vi sono: la sintesi di
proteine plasmatiche; la sintesi dei fattori della coagulazione; la
sintesi dell'ormone inattivo angiotensinogeno; la fagocitosi, tramite le
cellule di Kupffer, degli eritrociti danneggiati; l'accumulo di sangue,
principale contributo alle riserve venose dell'organismo;
l'assorbimento e la degradazione di ormoni circolanti (inclusi insulina
e adrenalina) e di immunoglobuline; l'assorbimento e l'inattivazione di
farmaci liposolubili.

Nella \textbf{regolazione metabolica} il fegato controlla i livelli di
lipidi, carboidrati e aminoacidi circolanti: tutto il sangue refluo
dalle superfici assorbenti dell'apparato digerente entra nel circolo
portale epatico e fluisce nel fegato, dove gli epatociti ne controllano i
livelli di metaboliti. Nella \textbf{regolazione ematologica} il fegato
rappresenta la principale riserva di sangue dell'organismo e riceve
circa il 25\% della gittata cardiaca; durante il passaggio del sangue
lungo i sinusoidi epatici, i fagociti rimuovono gli eritrociti
invecchiati o danneggiati, i microrganismi patogeni ed eventuali detriti
cellulari, mentre gli epatociti sintetizzano proteine plasmatiche,
contribuendo alla concentrazione osmotica del sangue, e trasportano
sostanze nutritizie. Nella \textbf{sintesi e secrezione di bile} gli
epatociti sintetizzano la bile, che viene accumulata nella cistifellea
ed escreta nel lume duodenale; la bile è costituita da acqua e ioni, che
diluiscono e neutralizzano gli acidi del chimo quando questo entra
nell'intestino tenue, e da sali biliari, che emulsionano i lipidi del
chimo.

% ---------------------------------------------------------------------------
\section{La cistifellea}

La cistifellea è un piccolo organo cavo piriforme, che può contenere
fino a 50--60 cm\textsuperscript{3} di bile. È situata nella fossa
cistica della faccia viscerale del fegato, trattenuta dal peritoneo che
la ricopre inferiormente e lateralmente; il suo collo ha un decorso
tortuoso che continua nel dotto cistico.

La funzione della cistifellea è il deposito e la modificazione della
bile: durante la permanenza nella cistifellea la composizione della bile
cambia, in quanto si verifica una concentrazione dei sali biliari. La
bile, sintetizzata dagli epatociti in ragione di circa 1 litro al giorno,
raggiunge la cistifellea, dove viene concentrata; da qui viene riversata
nel lume del duodeno, particolarmente dopo l'arrivo di grassi, in
seguito alla stimolazione delle cellule del sistema endocrino
gastroenteropatico a secernere l'ormone colecistochinina (CCK), che
svolge due compiti: stimolare la contrazione della muscolatura della
cistifellea e rilasciare il muscolo sfintere dell'ampolla, e stimolare il
pancreas a secernere nel duodeno gli enzimi digestivi che esso produce.
L'azione emulsionante della bile e quella idrolitica dei succhi
pancreatici ed enterici contribuiscono a completare la digestione dei
materiali nutritizi del chimo. L'eccessiva concentrazione dei sali
biliari nella cistifellea può favorire la precipitazione e la
formazione di calcoli biliari (colelitiasi).

% ---------------------------------------------------------------------------
\section{L'intestino tenue}

L'intestino tenue è deputato all'assorbimento e alla digestione delle
sostanze nutritizie: l'assorbimento avviene per il 90\% nell'intestino
tenue e solo per il 10\% nella porzione prossimale del crasso. È
suddiviso in una prima porzione, più breve (circa 25 cm) e meno mobile,
il \textbf{duodeno}, a forma di "C", e in una seconda porzione, più
lunga (circa 6,5 m) e più mobile, l'\textbf{intestino tenue
mesenteriale}, a sua volta suddiviso, senza un limite preciso, in
digiuno e ileo. Il peritoneo riveste l'intestino tenue mesenteriale e
solo la parte superiore del duodeno, mentre la parte discendente del
duodeno è rivestita dal peritoneo soltanto nei due terzi anteriori della
sua circonferenza; la diversa mobilità dei due tratti dipende proprio
dai rapporti con la tonaca sierosa.

L'intestino tenue mesenteriale è ripiegato in diverse anse e occupa
quasi interamente la cavità addominale sotto-mesocolica, incorniciato
superiormente e lateralmente dal colon; l'ileo termina confluendo nel
cieco, attraverso uno sfintere, la valvola ileocecale. La funzione di
questo tratto di intestino è, nella prima parte, di favorire la
degradazione dei singoli elementi che costituiscono il chimo
(carboidrati, proteine e lipidi) e, nella seconda parte, di provvedere al
loro riassorbimento.

% ---------------------------------------------------------------------------
\section{L'intestino crasso}

L'intestino crasso segue l'intestino tenue a livello della valvola
ileocecale e termina aprendosi all'esterno in corrispondenza dell'ano. È
lungo circa 1,6 m ed è situato prevalentemente nella cavità addominale,
tranne che per la sua parte terminale, il retto, che attraversa la
cavità pelvica e quindi il pavimento pelvico, per terminare all'esterno
nella regione perineale posteriore, nel canale anale.

All'intestino cieco è annessa l'appendice vermiforme. Il colon può
essere ulteriormente suddiviso in ascendente, trasverso, discendente e
sigmoideo. La maggior parte dell'intestino crasso ha ampi rapporti con
il peritoneo, anche se il colon ascendente, il colon discendente e il
retto ne sono rivestiti solo parzialmente. L'intestino retto è l'ultima
parte dell'intestino crasso e si suddivide in ampolla rettale, ospitata
nella pelvi, e canale anale, che attraversa il perineo e termina con
l'ano.

Le funzioni dell'intestino crasso sono: il riassorbimento di acqua ed
elettroliti e la compattazione del contenuto intestinale in feci;
l'assorbimento di vitamine liberate dalla flora batterica commensale;
l'accumulo del materiale fecale prima della defecazione. In sintesi, la
funzione principale dell'intestino crasso è favorire il riassorbimento
d'acqua per consolidare le feci.

% ---------------------------------------------------------------------------
\subsection{Nota clinica: patologie del colon-retto}

La proliferazione delle cellule della mucosa che rivestono internamente
la parete intestinale può dare origine a \textbf{polipi}, neoformazioni
carnose di forma variabile, che possono essere benigne o maligne. Tra le
neoplasie maligne del tratto di intestino crasso si annovera
l'\textbf{adenocarcinoma del colon-retto}, un tumore maligno che origina
dalla proliferazione incontrollata delle cellule della mucosa
intestinale.
